\clearpage
\VerdierPhysicalPageAnchor{0177}%
\begingroup\small\setlength{\parskip}{0pt}\setlength{\emergencystretch}{1.5em}
\setlength{\abovedisplayskip}{4pt plus 2pt minus 2pt}%
\setlength{\belowdisplayskip}{4pt plus 2pt minus 2pt}%

\VerdierSemanticAnchor{ch3d-1-2-1-truncations-cont}
\noindent（それぞれ
\[
\begin{aligned}
i^q&=\operatorname{id},&&q<n-1,\\
i^{n-1}&\colon\operatorname{Ker}(d_X^{n-1})\longrightarrow X^{n-1}
  &&\text{（標準射）},\\
i^q&=0,&&q\geq n
\end{aligned}
\]
）。

射 $p$（それぞれ $i$）は次数 $j\geq n$（それぞれ $j<n$）のコホモロジー
対象上で同型を誘導し、$j<n$ のとき $H^j(X(n,+\infty))=0$
（それぞれ $j\geq n$ のとき $H^j(X(-\infty,n))=0$）である。この構成を
用いれば、命題 (1.2.1) のすべての主張が直ちに証明される。詳細な証明は
読者に委ねる。

\VerdierSemanticAnchor{ch3d-1-2-2}
\medskip\noindent\textbf{定義 1.2.2.}\quad
三角圏
\[
D(\mathcal A)=K(\mathcal A)/\operatorname{Ac}(\mathcal A)
=K(\mathcal A)[\operatorname{Qis}(\mathcal A)^{-1}]
\]
を $\mathcal A$ の\emph{導来圏}と呼び、$D(\mathcal A)$ と書く
（第 II 章，2.2.6 および 2.2.10）。

定理 (2.2.6) により、標準コホモロジー的関手
$H\colon K(\mathcal A)\longrightarrow\mathcal A$ は、やはり同じ記号で表す
コホモロジー的関手を通じて一意に分解する：

\VerdierSemanticAnchor{ch3d-1-2-2-cohomology}
\begin{equation}
H\colon D(\mathcal A)\longrightarrow\mathcal A.
\tag{1.2.2.1}
\end{equation}
これも\emph{標準コホモロジー的関手}と呼ぶ。標準商関手を
$Q\colon K(\mathcal A)\longrightarrow D(\mathcal A)$ と書く。

\VerdierSemanticAnchor{ch3d-1-2-3}
\medskip\noindent\textbf{定理 1.2.3.}\quad
\textit{a) 次の図式に現れる自然関手は：}

\VerdierSemanticAnchor{ch3d-1-2-3-a-diagram}
\begin{center}
  \resizebox{0.80\textwidth}{!}{%
    $\begin{tikzcd}[ampersand replacement=\&,column sep=4.5em,row sep=3.0em]
      \& K^{,+}(\mathcal A)/\operatorname{Ac}(\mathcal A)
        \arrow[dr] \& \\
      K^{,b}(\mathcal A)/\operatorname{Ac}(\mathcal A)
        \arrow[ur] \arrow[dr]
        \& \& D(\mathcal A) \\
      \& K^{,-}(\mathcal A)/\operatorname{Ac}(\mathcal A)
        \arrow[ur] \&
    \end{tikzcd}$%
  }
\end{center}
\endgroup
